\begin{description}
\item[(a)] One can argue this in different ways. One can note that the
  Hubble parameter is often called Hubble constant, but actually it's not a
  constant. The dimensions of Hubble parameter is the inverse of time
  $[T^{-1}]$. Alternatively, one can simply look at unit of $H_0$ in the
  table of constants and conclude the same. \hfill(1 point)

  It's natural to define a timescale as the reciprocal of Hubble
  parameter: $t_H = \frac{1}{H}$, and this is the Hubble time which is a
  characteristic timescale of Universe expansion. \hfill(2 points)

  Present-day Hubble time $t_{H0} = 14.46$\,Gyr. \hfill(2 points)

\item[(b)] From Friedmann Equation, the critical density is defined in the
  way of
  $\left(\frac{\dot{a}}{a}\right)^2 = \frac{8\pi G}{3}\rho_c$, thus
  $\rho_c = \frac{3H^2}{8\pi G}$. \hfill(2 points)

  Substitute $H_0 = 2.19\times 10^{-18}\,\mathrm{s^{-1}}$ into above, we
  have $\rho_{c0} = 8.59\times 10^{-27}\ kg/m^3$ \hfill(3 points)

\item[(c)] $\Omega_{\mathrm{m}} + \Omega_{\mathrm{r}} + \Omega_\Lambda
  + \Omega_k = 1$
  \begin{align*}
  \frac{\rho_m}{\rho_c} + \frac{\rho_r}{\rho_c}
    + \Omega_\Lambda + \Omega_k &= 1\\
  \frac{8\pi G}{3H^2}(\rho_m + \rho_r)
    + \Omega_\Lambda + \Omega_k &= 1\\
  \frac{8\pi G}{3}(\rho_m + \rho_r)
    + H^2\Omega_\Lambda + H^2\Omega_k &= H^2
  \end{align*}
  \hfill(2 points)

  Comparing this with the Friedmann equation
  \[ H^2 = \frac{8\pi G}{3}(\rho_m + \rho_r)
     + \frac{\Lambda c^2}{3} - \frac{kc^2}{a^2}. \]
  \[ \Omega_\Lambda = \frac{\Lambda c^2}{3H^2}
     \quad\text{and}\quad
     \Omega_k = -\frac{kc^2}{a^2 H^2} \]
  \hfill(4 points)

\item[(d)] Pressure exerted by the radiation is
  $p = \frac{1}{3}\rho_r c^2$

  Thus, the Fluid Equation becomes,
  $\dot{\rho_r} + 3\frac{\dot{a}}{a}
   \left(\rho_r + \frac{1}{3}\rho_r\right) = 0$ \hfill(2 points)

  $\frac{\dot{\rho_r}}{\rho_r} = -4\frac{\dot{a}}{a}$ implying

  Thus $\rho_r \propto a^{-4}$. \hfill(3 points)

  And we know $a = \frac{1}{1+z}$, hence $\rho_r \propto (1+z)^4$.
  \hfill(2 points)

\item[(e)] In the fluid equation,
  $\dot{\rho} + 3\frac{\dot{a}}{a}(\rho + q\rho) = 0$.

  But the cosmological constant does not evolve. So
  $\dot{\rho_\Lambda} = 0$

  $\frac{3\dot{a}\rho_\Lambda}{a}(1 + q) = 0$ which means $q = -1$.
  \hfill(4 points)

\item[(f)] From the units, one can figure out that $\hbar$ has the
  dimensions of $[ML^2T^{-1}]$.

  $G$ has the dimensions of $[M^{-1}L^3T^{-2}]$.

  And speed of light $c$ has dimensions of $[LT^{-1}]$. \hfill(5 points)

  We assume that Planck time is given by $\hbar^x c^y G^z$. In order to
  make it a time dimensional quantity, the following equation must hold:
  $T = (ML^2T^{-2})^x(LT^{-1})^y(M^{-1}L^3T^{-2})^z$.
  % sic: source writes (ML^2 T^{-2})^x for hbar despite stating
  % [ML^2 T^{-1}] above; the exponent bookkeeping below uses T^{-1}
  $T = M^{x-z}L^{2x+y+3z}T^{-x-y-2z}$. \hfill(3 points)

  From $x - z = 0$, $2x + y + 3z = 0$, $-x - y - 2z = 1$, we get
  $x = \frac{1}{2}$, $y = \frac{-5}{2}$, $z = \frac{1}{2}$. Thus, the
  Planck time is approximately $t_p \approx \sqrt{\hbar G/c^5}$.
  \hfill(3 points)

  Substitute the numerical values into it and get:
  $t_P = 5.4\times 10^{-44}\,s$. \hfill(2 points)

\item[(g)] Schwarzschild radius of a black hole equals to $2l_P$:
  \[ r_s = \frac{2GM_P}{c^2} = 2ct_P = 2c\sqrt{\hbar G/c^5}, \]
  \hfill(3 points)
  \[ M_P = \sqrt{\hbar c/G}. \]
  \hfill(2 points)
  \[ M_P c^2 = \sqrt{\hbar c^5/G} = 1.22\times 10^{19}\,GeV. \]
  \hfill(2 points)

\item[(h)] Given the energy density, and the result of problem \textbf{d},
  we find that $\varepsilon_r \propto T^4$,
  $\varepsilon_r = \rho_r c^2$, $\rho_r \propto (1+z)^4$.

  Thus $T \propto 1 + z$, $\frac{T}{1+z} = const$. \hfill(4 points)

\item[(i)] The matter-radiation equality means
  $\Omega_r(z_{eq}) = \Omega_m(z_{eq})$. We need to derive the behavior of
  this two density parameters. Obviously
  $\Omega_m(z) = \Omega_{m0}(1+z)^3$ and
  $\Omega_r(z) = \Omega_{r0}(1+z)^4$. \hfill(4 points)

  Thus $1 + z_{eq} = \frac{\Omega_{m0}}{\Omega_{r0}}$. $\Omega_{r0}$ is
  the density parameter of CMB radiation and can be calculated.
  \hfill(1 point)

  By the definition of density parameter of photon radiation (footnote
  gamma: $\gamma$,
  \[ \Omega_{\gamma 0} = \frac{\rho_{\gamma 0}}{\rho_{c0}}
     = \frac{\varepsilon_{\gamma 0}}{c^2}\frac{8\pi G}{3H^2}
     = \frac{\varepsilon_{\gamma 0}}{c^2}
       \frac{8\pi G}{3H_{100}^2 h^2}, \]
  where $H_{100} = 100\ km/s/Mpc$. \hfill(4 points)

  The energy density of blackbody photon radiation is
  $\varepsilon_\gamma = \frac{\pi^2}{15\hbar^3 c^3}(k_B T)^4$, thus
  \[ \Omega_{\gamma 0}h^2 = \frac{\pi^2}{15\hbar^3 c^5}\,k_B^4\,
     \frac{8\pi G}{3H_{100}^2}\,T_0^4 = 2.47\times 10^{-5}. \]
  \hfill(4 points)

  But there is also an additional contribution from neutrinos of about
  68\%, so the total density parameter of radiation is
  $\Omega_{r0}h^2 = \Omega_{\gamma 0}h^2\times 1.68
   = 4.15\times 10^{-5}$. \hfill(2 points)

  Hence
  $1 + z_{eq} = \frac{\Omega_{m0}}{\Omega_{r0}}
   = 2.4\times 10^{4}\,\Omega_{m0}h^2$. \hfill(1 point)

\item[(j)] At that time, the redshift is
  \[ 1 + z = \frac{1\,MeV}{k_B\times 2.73\,K} = 4.25\times 10^{9}. \]

  Thus the scale factor
  $a = \frac{1}{1+z} = 2.35\times 10^{-10}$. \hfill(2 points)

  Only considering radiation dominated Universe, Friedmann Equation can be
  written as
  \[ \left(\frac{\dot{a}}{a}\right)^2 = \frac{8\pi G}{3}\rho_r
     = \frac{8\pi G\rho_{r0}}{3a^4}, \]
  \hfill(2 points)
  \[ \frac{\dot{a}}{a} = \frac{1}{a^2}\sqrt{\frac{8\pi G\rho_{r0}}{3}} \]
  \[ a^2 H = \sqrt{\frac{8\pi G\Omega_{r0}}{3}\rho_{c0}}
     = \sqrt{H_0^2\Omega_{r0}} \]
  \hfill(1 points)

  Thus, we have,
  \[ t = \frac{1}{2H} = \frac{a^2}{2\sqrt{H_0^2\Omega_{r0}}}
     = 1.3\,s \sim 1\,s. \]

  Thus neutrino decoupled at about 1\,s after big bang. \hfill(3 points)
\end{description}
